\chapter{Wasserstein 距離 \texorpdfstring{$W_p$}{Wp}}
\label{ch:wasserstein-metrics}

\section{設定と定義}
\label{sec:wp-definition}

$(X,d)$ を Polish 空間
（完備かつ可分な距離空間：定義~\ref{def:meas-polish}）とする．
Borel $\sigma$-代数 $\Borel(X)$ 上の確率測度全体を
\[
  \Prob(X)
  :=
  \left\{
    \mu\colon\Borel(X)\to[0,1]
    \;\middle|\;
    \mu\text{ は測度},\;\mu(X)=1
  \right\}
\]
とおく．

\begin{definition}{有限 $p$ 次モーメントをもつ確率測度}{wp-p-space}
  任意に点 $x_0\in X$ を一つ固定する．$1\leq p<\infty$ に対し
  \[
    \Pp{p}(X)
    :=
    \left\{
      \mu\in\Prob(X)
      \;\middle|\;
      \int_X d(x,x_0)^p\,\d\mu(x)<\infty
    \right\}
  \]
  と定める．また
  \begin{align*}
    \Pp{\infty}(X)
    &:=
    \left\{\mu\in\Prob(X)\mid \supp\mu\text{ は有界}\right\}\\
    &=
    \left\{
      \mu\in\Prob(X)
      \;\middle|\;
      \sup_{x\in\supp\mu}d(x,x_0)<\infty
    \right\}
  \end{align*}
  とおく（台 $\supp\mu$ は定義~\ref{def:meas-support}，
  有界集合は定義~\ref{def:meas-open-ball}を参照．
  2つの表示の一致は注意~\ref{rem:wp-p-infty-equiv}で確かめる）．
\end{definition}

\begin{lemma}{$\Pp{p}(X)$ の $x_0$ 非依存性}{wp-p-indep}
  $1\leq p<\infty$ のとき，$\Pp{p}(X)$ は点 $x_0$ の選び方によらない．
  （$p=\infty$ の場合，第1の表示（台の有界性）が $x_0$ に言及しないので，
  非依存性は明らかである．）
\end{lemma}

\begin{proof}
  $x_0,x_0'\in X$ と $\mu\in\Prob(X)$ を任意に取り，
  $\int_X d(x,x_0)^p\,\d\mu(x)<\infty$ を仮定する．
  三角不等式から
  \[
    d(x,x_0')
    \leq d(x,x_0)+d(x_0,x_0').
  \]
  両辺 $0$ 以上なので $p$ 乗して
  \[
    d(x,x_0')^p
    \leq
    \bigl(d(x,x_0)+d(x_0,x_0')\bigr)^p.
  \]
  ここで $a:=d(x,x_0)$，$b:=d(x_0,x_0')$ とおけば，
  $f(t)=t^p$（$p\geq 1$）が $[0,\infty)$ 上で凸であること
  （微積分の標準的な事実として認める）から
  \begin{align*}
    \text{右辺}
    = (a+b)^p
    &= 2^p\left(\frac{a}{2}+\frac{b}{2}\right)^p\\
    &\leq 2^p\left(\frac{a^p}{2}+\frac{b^p}{2}\right)
    = 2^{p-1}(a^p+b^p).
  \end{align*}
  よって
  \[
    d(x,x_0')^p
    \leq
    2^{p-1}\bigl(d(x,x_0)^p+d(x_0,x_0')^p\bigr).
  \]
  両辺を $\d\mu(x)$ で積分すると
  （右辺第2項は定数なので，$\mu(X)=1$ より積分値は $d(x_0,x_0')^p$）
  \[
    \int_X d(x,x_0')^p\,\d\mu(x)
    \leq
    2^{p-1}\left(\int_X d(x,x_0)^p\,\d\mu(x)+d(x_0,x_0')^p\right)
    <\infty.
  \]
  $x_0$ と $x_0'$ を入れ替えれば逆も成り立つ．
\end{proof}

\begin{remark}{$\Pp{\infty}(X)$ の言い換え}{wp-p-infty-equiv}
  定義~\ref{def:wp-p-space}の $\Pp{\infty}(X)$ の2つの表示が
  一致することを確かめる．
  \begin{itemize}
    \item $\supp\mu$ が有界，すなわちある開球 $B_r(z)$ に含まれるとする．
      三角不等式より $\supp\mu\subset B_{r+d(z,x_0)}(x_0)$ だから，
      $R:=r+d(z,x_0)$ とおけば，すべての $x\in\supp\mu$ で
      $d(x,x_0)<R$，よって
      $\sup_{x\in\supp\mu}d(x,x_0)\leq R<\infty$ である．
    \item 逆に $S:=\sup_{x\in\supp\mu}d(x,x_0)<\infty$ とすると，
      すべての $x\in\supp\mu$ で $d(x,x_0)\leq S<S+1$，
      すなわち $\supp\mu\subset B_{S+1}(x_0)$ だから
      $\supp\mu$ は有界である．
  \end{itemize}
\end{remark}

\begin{definition}{Coupling 集合}{wp-coupling}
  $\mu,\nu\in\Prob(X)$ に対して，\textbf{Coupling 集合}を
  \[
    \Couplings(\mu,\nu)
    :=
    \left\{
      \pi\in\Prob(X\times X)
      \;\middle|\;
      \begin{aligned}
        \pi(A\times X)&=\mu(A) &&(\forall A\in\Borel(X)),\\
        \pi(X\times B)&=\nu(B) &&(\forall B\in\Borel(X))
      \end{aligned}
    \right\}
  \]
  と定める．
  この条件は「$\pi$ の第1周辺分布が $\mu$，第2周辺分布が $\nu$」，
  すなわち射影による押し出しで
  ${\mathrm{pr}_1}\pushforward\pi=\mu$，${\mathrm{pr}_2}\pushforward\pi=\nu$
  と言い換えられる（補題~\ref{lem:meas-marginal-integral}）．
\end{definition}

\begin{remark}{記法：$(x,y)\sim\pi$ と a.e.}{wp-sampling-notation}
  $\pi\in\Prob(X\times X)$ に対し，$(x,y)\sim\pi$ は
  「確率測度 $\pi$ に従って選ばれた点対 $(x,y)$」を表す文献の記法で，
  $\int_{X\times X} f(x,y)\,\d\pi(x,y)$ は
  $\mathbb{E}_{(x,y)\sim\pi}[f(x,y)]$ とも書かれる
  （本資料では積分記法を用いる）．
  性質が $\pi$-a.e. で成り立つ，の意味は
  定義~\ref{def:meas-ae}の通りである．
\end{remark}

\begin{definition}{Wasserstein 距離}{wp-wasserstein}
  $1\leq p<\infty$ と $\mu,\nu\in\Pp{p}(X)$ に対し，
  下限（定義~\ref{def:meas-sup-inf}）を用いて
  \[
    \Wass_p(\mu,\nu)
    :=
    \inf_{\pi\in\Couplings(\mu,\nu)}
    \left(
      \int_{X\times X}d(x,y)^p\,\d\pi(x,y)
    \right)^{1/p},
  \]
  また $\mu,\nu\in\Pp{\infty}(X)$ に対し，
  $L^\infty$ ノルム（定義~\ref{def:meas-lp-norm}）を用いて
  \[
    \Wass_\infty(\mu,\nu)
    :=
    \inf_{\pi\in\Couplings(\mu,\nu)}
    \norm{d}_{L^\infty(\pi)}
  \]
  と定める．
\end{definition}

\begin{remark}{記法と可測性}{wp-welldef}
  距離関数 $d\colon X\times X\to[0,\infty)$ は，三角不等式から従う評価
  \[
    \abs{d(x,y)-d(x',y')}\leq d(x,x')+d(y,y')
  \]
  により連続，したがって Borel 可測である
  （注意~\ref{rem:meas-measurable-ops}）．
  $d\geq0$ より $1\leq p<\infty$ では
  $\norm{d}_{L^p(\pi)}=\bigl(\int d^p\,\d\pi\bigr)^{1/p}$ だから，
  $1\leq p\leq\infty$ を通じて
  \[
    \Wass_p(\mu,\nu)
    =\inf_{\pi\in\Couplings(\mu,\nu)}\norm{d}_{L^p(\pi)}
  \]
  と統一的に書ける．
  また主張~\ref{clm:meas-esssup-ae}より，$\norm{d}_{L^\infty(\pi)}$ 自身が
  「$d(x,y)\leq R$ が $\pi$-a.e. で成り立つ」をみたす最小の $R\geq0$ であり，
  これは $\esssup_{(x,y)\sim\pi}d(x,y)$ とも書かれる．
\end{remark}

\begin{lemma}{$\Wass_p$ の有限性}{wp-finite}
  $1\leq p\leq\infty$，$\mu,\nu\in\Pp{p}(X)$ とすると，
  $\Couplings(\mu,\nu)\neq\emptyset$ であり，
  $\Wass_p(\mu,\nu)<\infty$ である．
\end{lemma}

\begin{proof}
  直積測度 $\mu\otimes\nu$（定理~\ref{thm:meas-product}）は，
  主張~\ref{clm:meas-product-marginals}より
  第1周辺分布 $\mu$・第2周辺分布 $\nu$ をもつから，
  $\mu\otimes\nu\in\Couplings(\mu,\nu)$，
  特に $\Couplings(\mu,\nu)\neq\emptyset$ である．
  下限の定義から
  $\Wass_p(\mu,\nu)\leq\norm{d}_{L^p(\mu\otimes\nu)}$
  なので，右辺が有限であることを示せばよい．

  $1\leq p<\infty$ のとき．
  補題~\ref{lem:wp-p-indep}の証明と同じ凸性の評価から，各点で
  \[
    d(x,y)^p
    \leq\bigl(d(x,x_0)+d(x_0,y)\bigr)^p
    \leq2^{p-1}\bigl(d(x,x_0)^p+d(y,x_0)^p\bigr)
  \]
  が成り立つ．周辺積分公式（補題~\ref{lem:meas-marginal-integral}）より
  \[
    \int_{X\times X}d(x,x_0)^p\,\d(\mu\otimes\nu)
    =\int_X d(x,x_0)^p\,\d\mu(x),
    \qquad
    \int_{X\times X}d(y,x_0)^p\,\d(\mu\otimes\nu)
    =\int_X d(y,x_0)^p\,\d\nu(y)
  \]
  だから，積分の単調性と合わせて
  \[
    \int_{X\times X}d(x,y)^p\,\d(\mu\otimes\nu)
    \leq
    2^{p-1}\left(
      \int_X d(x,x_0)^p\,\d\mu+\int_X d(y,x_0)^p\,\d\nu
    \right)
    <\infty.
  \]

  $p=\infty$ のとき．
  $R_\mu:=\sup_{x\in\supp\mu}d(x,x_0)$，
  $R_\nu:=\sup_{y\in\supp\nu}d(y,x_0)$ とおくと，
  $\mu,\nu\in\Pp{\infty}(X)$ よりいずれも有限である．
  $X$ は可分だから台は全測度をもち（補題~\ref{lem:meas-support-full}），
  \[
    (\mu\otimes\nu)(\supp\mu\times\supp\nu)
    =\mu(\supp\mu)\,\nu(\supp\nu)=1.
  \]
  $\supp\mu\times\supp\nu$ 上では三角不等式より
  $d(x,y)\leq d(x,x_0)+d(x_0,y)\leq R_\mu+R_\nu$ だから，
  「$d\leq R_\mu+R_\nu$」は $(\mu\otimes\nu)$-a.e. で成り立つ．
  よって $\norm{d}_{L^\infty(\mu\otimes\nu)}\leq R_\mu+R_\nu<\infty$ である．
\end{proof}

\section{Wasserstein 距離の距離性}
\label{sec:metric-property}

\begin{definition}{Lipschitz 関数}{wp-lipschitz-function}
  距離空間 $(X,d)$ 上の関数 $f\colon X\to\R$ が
  \textbf{Lipschitz 関数}であるとは，ある定数 $L\geq0$ が存在して
  \[
    \abs{f(x)-f(y)}\leq Ld(x,y)
    \qquad(\forall x,y\in X)
  \]
  が成り立つことをいう．
  この条件をみたす $L$ 全体の下限 $L^*$ もまた条件をみたす
  （条件をみたす列 $L_n\downarrow L^*$ を取り，
  各 $x,y$ で両辺の極限を取ればよい）ので，$L^*$ は最小値である．
  これを $f$ の \textbf{Lipschitz 定数}とよび $\Lip(f)$ と書く．
  さらに上限（定義~\ref{def:meas-sup-inf}）の意味で
  $\sup_{x\in X}\abs{f(x)}<\infty$ であるとき，
  $f$ は\textbf{有界 Lipschitz 関数}であるという．
\end{definition}

\begin{lemma}{有界 Lipschitz 関数による測度の判定}{wp-bl-separates}
  $\mu,\nu\in\Prob(X)$ とする．すべての有界 Lipschitz 関数 $f\colon X\to\R$ について
  \[
    \int_X f\,\d\mu=\int_X f\,\d\nu
  \]
  が成り立つならば，$\mu=\nu$ である．
\end{lemma}

\begin{proof}
  方針：閉集合の指示関数を Lipschitz 関数で上から近似し，
  有界収束定理で「すべての閉集合上で $\mu=\nu$」に落とす．

  閉集合 $F\subset X$ を任意に取る．
  $F=\emptyset$ のときは $\mu(\emptyset)=0=\nu(\emptyset)$ だから，
  以下 $F\neq\emptyset$ とする．
  $d(x, F) := \inf_{y \in F} d(x, y)$ とおき，各 $n \in \N$ で
  \[
    f_n(x) := \max\{1 - n \, d(x, F),\; 0\}
  \]
  を定める．
  $f_n$ は $0 \leq f_n \leq 1$ をみたす有界 Lipschitz 関数である．実際，
  任意の $z\in F$ で $d(x,F)\leq d(x,z)\leq d(x,y)+d(y,z)$ だから
  $z$ について下限を取れば $d(x,F)\leq d(x,y)+d(y,F)$，
  $x$ と $y$ を入れ替えて
  \[
    \abs{d(x,F)-d(y,F)}\leq d(x,y)
  \]
  であり，さらに $s,t\geq0$ に対し
  $\abs{\max\{1-ns,0\}-\max\{1-nt,0\}}\leq n\abs{s-t}$ だから
  \[
    \abs{f_n(x)-f_n(y)}
    \leq n\,\abs{d(x,F)-d(y,F)}
    \leq n\,d(x,y).
  \]
  したがって仮定より，各 $n$ について
  \[
    \int_X f_n \, \d\mu = \int_X f_n \, \d\nu
  \]
  が成り立つ．

  ここで $F$ は空でない閉集合だから，
  $d(x, F) = 0 \Leftrightarrow x \in F$ である．
  \begin{itemize}
    \item $x \in F$ のとき：$d(x, F) = 0$ より $f_n(x) = 1$．
    \item $x \notin F$ のとき：$d(x, F) > 0$ より，
      十分大きい $n$ で $n \, d(x, F) \geq 1$ となり $f_n(x) = 0$．
  \end{itemize}
  よって $f_n(x) \to \mathbf{1}_F(x)$（すべての $x \in X$）である．
  さらに $0 \leq f_n \leq 1$，$f_n$ は連続ゆえ可測，
  $\mu$ は有限測度（確率測度）だから，
  有界収束定理（定理~\ref{thm:meas-bounded-convergence}）により
  \[
    \mu(F)
    = \int_X \mathbf{1}_F \, \d\mu
    = \lim_{n \to \infty} \int_X f_n \, \d\mu
    = \lim_{n \to \infty} \int_X f_n \, \d\nu
    = \int_X \mathbf{1}_F \, \d\nu
    = \nu(F).
  \]
  $F$ は任意の閉集合だったから，$\mu$ と $\nu$ はすべての閉集合上で一致する．
  閉集合の全体 $\mathcal{C}$ は，$X\in\mathcal{C}$ であり，
  2つの閉集合の共通部分が閉集合であることから
  $\pi$-系である（定理~\ref{thm:meas-uniqueness}）．
  開集合は閉集合の補集合だから，
  $\mathcal{C}$ が生成する $\sigma$-代数は $\Borel(X)$ に一致する．
  $\mu(X)=\nu(X)=1$ だから，
  一致定理（定理~\ref{thm:meas-uniqueness}）より $\mu = \nu$ である．
\end{proof}

\begin{lemma}{三変数 coupling の存在（gluing lemma）}{wp-gluing}
  $\mu,\nu,\xi\in\Prob(X)$，$\pi\in\Couplings(\mu,\nu)$，
  $\sigma\in\Couplings(\nu,\xi)$ とする．
  このとき $X^3$ 上の確率測度 $\gamma$ で，
  $(x,y)$-周辺分布が $\pi$，$(y,z)$-周辺分布が $\sigma$
  であるものが存在する．
  ここで $(x,y)$-周辺分布とは，射影
  $\mathrm{pr}_{12}\colon(x,y,z)\mapsto(x,y)$ による押し出し
  ${\mathrm{pr}_{12}}\pushforward\gamma$ のことをいう（$(y,z)$ も同様）．
\end{lemma}

\begin{proof}
  方針：$\pi$ と $\sigma$ がともにもつ周辺分布 $\nu$ を軸に，
  $y$ で条件付けた2つの Markov 核を貼り合わせる．

  $X$ は Polish 空間だから，
  分解定理（定理~\ref{thm:meas-disintegration}）が使える．
  $\pi$ の第2周辺分布は $\nu$ だから，Markov 核 $K_\pi$ で
  \[
    \pi(A\times B)=\int_B K_\pi(y,A)\,\d\nu(y)
    \qquad(A,B\in\Borel(X))
  \]
  をみたすものが取れる．
  次に $s(y,z):=(z,y)$ とおくと（連続ゆえ可測），
  $(s\pushforward\sigma)(C\times B)
  =\sigma(s^{-1}(C\times B))=\sigma(B\times C)$ だから，
  $s\pushforward\sigma$ の第2周辺分布は $\sigma$ の第1周辺分布，
  すなわち $\nu$ である．
  再び分解定理を $s\pushforward\sigma$ に適用して，Markov 核 $K_\sigma$ で
  \[
    \sigma(B\times C)
    =(s\pushforward\sigma)(C\times B)
    =\int_B K_\sigma(y,C)\,\d\nu(y)
    \qquad(B,C\in\Borel(X))
  \]
  をみたすものが取れる．

  定理~\ref{thm:meas-kernel-composition}により，$X^3$ 上の確率測度
  \[
    \gamma(\d x,\d y,\d z)
    :=
    K_\pi(y,\d x)\,\nu(\d y)\,K_\sigma(y,\d z),
  \]
  すなわち
  $\gamma(A\times B\times C)=\int_B K_\pi(y,A)\,K_\sigma(y,C)\,\d\nu(y)$
  をみたす $\gamma$ が存在する．
  2つの周辺分布を検証する．
  ${\mathrm{pr}_{12}}^{-1}(A\times B)=A\times B\times X$ と
  $K_\sigma(y,X)=1$（Markov 核の値は確率測度）から，
  任意の $A,B\in\Borel(X)$ に対し
  \[
    ({\mathrm{pr}_{12}}\pushforward\gamma)(A\times B)
    =\gamma(A\times B\times X)
    =\int_B K_\pi(y,A)\,K_\sigma(y,X)\,\d\nu(y)
    =\int_B K_\pi(y,A)\,\d\nu(y)
    =\pi(A\times B).
  \]
  矩形集合の全体は $\Borel(X\times X)$ を生成する $\pi$-系だから
  （注意~\ref{rem:meas-borel-product}），
  一致定理（定理~\ref{thm:meas-uniqueness}）より
  ${\mathrm{pr}_{12}}\pushforward\gamma=\pi$ である．
  同様に $K_\pi(y,X)=1$ より
  \[
    ({\mathrm{pr}_{23}}\pushforward\gamma)(B\times C)
    =\gamma(X\times B\times C)
    =\int_B K_\pi(y,X)\,K_\sigma(y,C)\,\d\nu(y)
    =\int_B K_\sigma(y,C)\,\d\nu(y)
    =\sigma(B\times C)
  \]
  だから ${\mathrm{pr}_{23}}\pushforward\gamma=\sigma$ である．
\end{proof}

\begin{proposition}{Wasserstein 距離は距離である}{wp-metric}
  $1\leq p\leq\infty$ に対し，
  $\Wass_p$ は $\Pp{p}(X)$ 上の距離である．
\end{proposition}

\begin{proof}
  距離の公理（定義~\ref{def:meas-metric-space}）を順に確かめる．
  以下，統一記法
  $\Wass_p(\mu,\nu)=\inf_{\pi}\norm{d}_{L^p(\pi)}$
  （注意~\ref{rem:wp-welldef}）を $1\leq p\leq\infty$ で共通に用いる．

  \textbf{(0) 非負性と有限性．}
  任意の $\pi\in\Couplings(\mu,\nu)$ で $\norm{d}_{L^p(\pi)}\geq0$ だから
  $\Wass_p(\mu,\nu)\geq0$，
  また補題~\ref{lem:wp-finite}より
  $\Couplings(\mu,\nu)\neq\emptyset$ かつ $\Wass_p(\mu,\nu)<\infty$．
  よって $\Wass_p$ は $\Pp{p}(X)$ 上の $[0,\infty)$ に値をもつ関数である．

  \textbf{(1) 対称性．}
  $s(x,y):=(y,x)$ とおく（連続ゆえ可測）．
  $\pi\in\Couplings(\mu,\nu)$ に対し $\pi':=s\pushforward\pi$ とおくと，
  \[
    \pi'(B\times X)=\pi(s^{-1}(B\times X))=\pi(X\times B)=\nu(B),
    \qquad
    \pi'(X\times A)=\pi(A\times X)=\mu(A)
  \]
  だから $\pi'\in\Couplings(\nu,\mu)$ である．
  コストは，系~\ref{cor:meas-pf-lp}と $d$ の対称性（$d\circ s=d$）より
  \[
    \norm{d}_{L^p(\pi')}
    =\norm{d\circ s}_{L^p(\pi)}
    =\norm{d}_{L^p(\pi)}.
  \]
  $s\circ s=\Id$ と合成則（主張~\ref{clm:meas-pf-compose}）より，
  対応 $\pi\mapsto s\pushforward\pi$ は
  $\Couplings(\mu,\nu)$ から $\Couplings(\nu,\mu)$ への全単射だから，
  両辺の下限を取って $\Wass_p(\mu,\nu)=\Wass_p(\nu,\mu)$ である．

  \textbf{(2) 一致の公理．}
  まず $\Wass_p(\mu,\mu)=0$ を示す．
  対角写像 $\iota(x):=(x,x)$（連続ゆえ可測）による押し出し
  $\pi^*:=\iota\pushforward\mu$ は，$\iota^{-1}(A\times B)=A\cap B$ より
  \[
    \pi^*(A\times X)=\mu(A\cap X)=\mu(A),
    \qquad
    \pi^*(X\times B)=\mu(B)
  \]
  をみたすから $\pi^*\in\Couplings(\mu,\mu)$ である．
  $d\circ\iota\equiv0$ だから，系~\ref{cor:meas-pf-lp}より
  $1\leq p\leq\infty$ で一括して
  \[
    \norm{d}_{L^p(\pi^*)}=\norm{d\circ\iota}_{L^p(\mu)}=0,
  \]
  よって $\Wass_p(\mu,\mu)=0$ である．

  逆に $\Wass_p(\mu,\nu)=0$ とする．
  任意の $\varepsilon>0$ に対し，下限の性質から
  $\pi_\varepsilon\in\Couplings(\mu,\nu)$ で
  \[
    \norm{d}_{L^p(\pi_\varepsilon)}<\varepsilon
  \]
  となるものを取れる．
  有界 Lipschitz 関数 $f\colon X\to\R$ を任意に取る．
  周辺積分公式（補題~\ref{lem:meas-marginal-integral}）より
  $\int_X f\,\d\mu=\int f(x)\,\d\pi_\varepsilon$，
  $\int_X f\,\d\nu=\int f(y)\,\d\pi_\varepsilon$ だから，
  $\abs{f(x)-f(y)}\leq\Lip(f)\,d(x,y)$ と合わせて
  \[
    \left|\int_X f\,\d\mu-\int_X f\,\d\nu\right|
    =
    \left|\int_{X\times X}(f(x)-f(y))\,\d\pi_\varepsilon\right|
    \leq
    \Lip(f)\int_{X\times X}d(x,y)\,\d\pi_\varepsilon
    =
    \Lip(f)\,\norm{d}_{L^1(\pi_\varepsilon)}.
  \]
  確率測度上では $\norm{d}_{L^1}\leq\norm{d}_{L^p}$
  （主張~\ref{clm:meas-l1-lp}，$p=\infty$ を含む）だから
  \[
    \left|\int_X f\,\d\mu-\int_X f\,\d\nu\right|
    \leq
    \Lip(f)\,\norm{d}_{L^p(\pi_\varepsilon)}
    \leq
    \Lip(f)\,\varepsilon.
  \]
  $\varepsilon>0$ は任意だったから，すべての有界 Lipschitz 関数 $f$ で
  $\int f\,\d\mu=\int f\,\d\nu$ となり，
  補題~\ref{lem:wp-bl-separates}より $\mu=\nu$ である．

  \textbf{(3) 三角不等式．}
  方針：ほぼ最適な $\pi,\sigma$ を gluing lemma で貼り合わせ，
  $(x,z)$-周辺分布を $\mu$ から $\xi$ への coupling として使う．

  $\mu,\nu,\xi\in\Pp{p}(X)$ と $\varepsilon>0$ を任意に取る．
  下限の性質から，$\pi\in\Couplings(\mu,\nu)$ と
  $\sigma\in\Couplings(\nu,\xi)$ で
  \[
    \norm{d}_{L^p(\pi)}<\Wass_p(\mu,\nu)+\varepsilon,
    \qquad
    \norm{d}_{L^p(\sigma)}<\Wass_p(\nu,\xi)+\varepsilon
  \]
  なるものを取れる（(0) よりいずれの $\Wass_p$ も有限であることに注意）．
  補題~\ref{lem:wp-gluing}により，$X^3$ 上の確率測度 $\gamma$ で
  ${\mathrm{pr}_{12}}\pushforward\gamma=\pi$，
  ${\mathrm{pr}_{23}}\pushforward\gamma=\sigma$
  なるものが取れる．
  $\tau:={\mathrm{pr}_{13}}\pushforward\gamma$ とおくと，
  \[
    \tau(A\times X)=\gamma(A\times X\times X)=\pi(A\times X)=\mu(A),
    \qquad
    \tau(X\times C)=\gamma(X\times X\times C)=\sigma(X\times C)=\xi(C)
  \]
  だから $\tau\in\Couplings(\mu,\xi)$ であり，
  \[
    \Wass_p(\mu,\xi)
    \leq \norm{d}_{L^p(\tau)}.
  \]
  以下，$\gamma$ 上の $L^p$ ノルムでは，
  $d(x,z)$ 等を $(x,y,z)\mapsto d(x,z)$ の略記とする．
  系~\ref{cor:meas-pf-lp}（$T=\mathrm{pr}_{13}$）より
  \[
    \norm{d}_{L^p(\tau)}=\norm{d(x,z)}_{L^p(\gamma)}
  \]
  であり，距離の三角不等式 $d(x,z)\leq d(x,y)+d(y,z)$（各点）と
  $L^p$ ノルムの単調性（主張~\ref{clm:meas-lp-monotone}），
  続いて Minkowski の不等式（定理~\ref{thm:meas-minkowski}）から
  \[
    \norm{d(x,z)}_{L^p(\gamma)}
    \leq \norm{d(x,y)+d(y,z)}_{L^p(\gamma)}
    \leq \norm{d(x,y)}_{L^p(\gamma)}+\norm{d(y,z)}_{L^p(\gamma)}.
  \]
  再び系~\ref{cor:meas-pf-lp}（$T=\mathrm{pr}_{12}$ と $T=\mathrm{pr}_{23}$）より
  $\norm{d(x,y)}_{L^p(\gamma)}=\norm{d}_{L^p(\pi)}$，
  $\norm{d(y,z)}_{L^p(\gamma)}=\norm{d}_{L^p(\sigma)}$ だから，まとめて
  \[
    \Wass_p(\mu,\xi)
    \leq\norm{d}_{L^p(\pi)}+\norm{d}_{L^p(\sigma)}
    <\Wass_p(\mu,\nu)+\Wass_p(\nu,\xi)+2\varepsilon.
  \]
  $\varepsilon>0$ は任意だったから
  \[
    \Wass_p(\mu,\xi)
    \leq
    \Wass_p(\mu,\nu)+\Wass_p(\nu,\xi)
  \]
  が従う．

  以上 (0)--(3) より，$\Wass_p$ は $\Pp{p}(X)$ 上の距離である．
\end{proof}
